% GENERATED FILE. Edit canonical lessons, then run npm run generate:books.
% canonical-source-hash: fnv1a64:39148e4d95ad4481
% canonical-lessons: SA-C262-asman, SA-C262-yusman, SA-C262-object-person-number, SA-C262-object-pronoun-retrieval

\chapter{Us and you all as objects}
\label{ch:sa-us-and-you-all-objects}
\hlchaptermodality{262}

\begin{chapteropening}
\textbf{By the end of this chapter:} \emph{I can choose first or second person and one, two, or several people for a known object frame.}

Complete chapter 262 by retrieving all six independent first- and second-person objects from person-and-number cues.
\end{chapteropening}

\section[\texorpdfstring{\sk{अस्मान्} (asmān)}{अस्मान् (asmān)}]{\sk{अस्मान्} (asmān) — us, several people, as the object}

\label{lesson:SA-C262-asman}



\begin{quote}
“Us two” is \textbf{\sk{आवाम्}} (\emph{āvām}). “Our” for several possessors is \textbf{\sk{अस्माकम्}} (\emph{asmākam}). Keep the \textbf{asm-} opening and make several people the object.
\end{quote}

\subsection*{You'll want to know: \sk{अस्मान्}}
\textbf{\sk{अस्मान्}} (\emph{asmān}) means “us” when several people including the speaker are the object.

\begin{quote}
\textbf{\sk{सः} \sk{अस्मान्} \sk{पश्यति।}}
\end{quote}

\begin{quote}
\emph{saḥ asmān paśyati.}
\end{quote}

\begin{quote}
\textbf{He sees us.}
\end{quote}

Whitney §491 lists \textbf{\sk{अस्मान्}} as the independent accusative plural. Do not confuse it with \textbf{\sk{अस्माकम्}}, “our / of us”: \textbf{asmān} is a person reached by the action; \textbf{asmākam} answers “whose?”.

\begin{cousinweb}
As a grammatical form, \textbf{\sk{अस्मान्}} comes from the first-person pronoun paradigm; it is not a word from a separate lexical root. Whitney's \emph{Sanskrit Grammar} §491 records it as the independent accusative plural. This lesson makes no deeper historical derivation claim.

Source: Whitney, \emph{Sanskrit Grammar} (1889), §491.
\end{cousinweb}

\subsection*{Your turn}
\begin{itemize}
\raggedright
  \item \emph{Say it:} “us two” as object — \emph{āvām}
  \item \emph{Say it:} “us, several people” as object — \emph{asmān}
  \item \emph{Say it:} “he sees us” — \emph{saḥ asmān paśyati}
  \item \emph{Contrast:} object \emph{asmān}; possessor \emph{asmākam}
\end{itemize}

\begin{tcolorbox}[breakable,colback=teal!4,colframe=teal!35!black,title={Before you move on}]
Several people including the speaker are seen. Which object form means “us”? (\textbf{\sk{अस्मान्}}, \emph{asmān}.)
\end{tcolorbox}
\section[\texorpdfstring{\sk{युष्मान्} (yuṣmān)}{युष्मान् (yuṣmān)}]{\sk{युष्मान्} (yuṣmān) — you, several people, as the object}

\label{lesson:SA-C262-yusman}



\begin{quote}
“Us, several people” is \textbf{\sk{अस्मान्}} (\emph{asmān}). “Your” to several people is \textbf{\sk{युष्माकम्}} (\emph{yuṣmākam}). Keep the \textbf{yuṣm-} opening and make the several addressees the object.
\end{quote}

\subsection*{You'll want to know: \sk{युष्मान्}}
\textbf{\sk{युष्मान्}} (\emph{yuṣmān}) means “you” when several familiar addressees are the object.

\begin{quote}
\textbf{\sk{सः} \sk{युष्मान्} \sk{पश्यति।}}
\end{quote}

\begin{quote}
\emph{saḥ yuṣmān paśyati.}
\end{quote}

\begin{quote}
\textbf{He sees you all.}
\end{quote}

Whitney §491 lists \textbf{\sk{युष्मान्}} as the independent accusative plural. Keep it apart from \textbf{\sk{युष्माकम्}}, “your / of you all”: the first is an object and the second is a possessor.

\begin{cousinweb}
As a grammatical form, \textbf{\sk{युष्मान्}} comes from the second-person pronoun paradigm; it is not a word from a separate lexical root. Whitney's \emph{Sanskrit Grammar} §491 records it as the independent accusative plural. This lesson makes no deeper historical derivation claim.

Source: Whitney, \emph{Sanskrit Grammar} (1889), §491.
\end{cousinweb}

\subsection*{Your turn}
\begin{itemize}
\raggedright
  \item \emph{Say it:} “he sees us” — \emph{saḥ asmān paśyati}
  \item \emph{Say it:} “he sees you all” — \emph{saḥ yuṣmān paśyati}
  \item \emph{Say it:} several people with the speaker — \emph{asmān}; several addressees — \emph{yuṣmān}
  \item \emph{Contrast:} object \emph{yuṣmān}; possessor \emph{yuṣmākam}
\end{itemize}

\begin{tcolorbox}[breakable,colback=teal!4,colframe=teal!35!black,title={Before you move on}]
Several familiar addressees are seen. Which object form means “you all”? (\textbf{\sk{युष्मान्}}, \emph{yuṣmān}.)
\end{tcolorbox}
\section[(object person and number)]{Who is seen, and how many?}

\label{lesson:SA-C262-object-person-number}



\begin{quote}
Give the singular pair, then the dual pair, then the several-person pair. (\emph{Mām, tvām; āvām, yuvām; asmān, yuṣmān.})
\end{quote}

\begin{grammarlens}[title={Grammar Lens: person first, then number}]
Choose in two small steps.

First ask who the object is:

\begin{itemize}
\raggedright
  \item first person includes the speaker;
  \item second person is the familiar person or people addressed.
\end{itemize}

Then count the people:

\begin{itemize}
\raggedright
  \item one person: \textbf{\sk{माम्}} / \textbf{\sk{त्वाम्}} (\emph{mām / tvām});
  \item two people: \textbf{\sk{आवाम्}} / \textbf{\sk{युवाम्}} (\emph{āvām / yuvām});
  \item several people: \textbf{\sk{अस्मान्}} / \textbf{\sk{युष्मान्}} (\emph{asmān / yuṣmān}).
\end{itemize}

The six are independent full forms. Choose person, choose number, then place the chosen object before the known verb.
\end{grammarlens}

\subsection*{Your turn}
Say these aloud:

\begin{itemize}
\raggedright
  \item first person, one — \emph{mām}; first person, two — \emph{āvām}; first person, several — \emph{asmān}
  \item second person, one — \emph{tvām}; second person, two — \emph{yuvām}; second person, several — \emph{yuṣmān}
  \item “he sees us two” — \emph{saḥ āvām paśyati}
  \item “he sees you all” — \emph{saḥ yuṣmān paśyati}
\end{itemize}

\begin{tcolorbox}[breakable,colback=teal!4,colframe=teal!35!black,title={Before you move on}]
What are the two choices you make before selecting the object form? (Person: speaker included or addressee; then number: one, two, or several.)
\end{tcolorbox}
\section[Practice]{One, two, or several — me/us or you}

\label{lesson:SA-C262-object-pronoun-retrieval}



\begin{quote}
Say the decision aloud: person first, then number. Do not recite a table; answer each scene separately.
\end{quote}

\subsection*{Your turn — choose before you speak}
Say these aloud:

\begin{itemize}
\raggedright
  \item the speaker alone is seen — \emph{saḥ mām paśyati}
  \item one familiar addressee is seen — \emph{saḥ tvām paśyati}
  \item the speaker and one companion are seen — \emph{saḥ āvām paśyati}
  \item two familiar addressees are seen — \emph{saḥ yuvām paśyati}
  \item several people including the speaker are seen — \emph{saḥ asmān paśyati}
  \item several familiar addressees are seen — \emph{saḥ yuṣmān paśyati}
\end{itemize}

\begin{grammarlens}[title={Grammar Lens: the boundary of this chapter}]
These six are the independent full forms taught here. Classical Sanskrit also has short alternatives in particular positions, and it has third-person and respectful object forms. Those belong to later lessons. This set is \textbf{not the whole Classical Sanskrit pronoun system}; it is the complete first- and second-person full-form choice across one person, two people, and several people.

Source for the full and short forms: William Dwight Whitney, \emph{Sanskrit Grammar} (1889), §491.
\end{grammarlens}

\begin{tcolorbox}[breakable,colback=teal!4,colframe=teal!35!black,title={Before you move on}]
Without a table, answer six prompts: me, you one, us two, you two, us several, you several. (\emph{Mām; tvām; āvām; yuvām; asmān; yuṣmān.})
\end{tcolorbox}
